\section{Момент инерции}

\subsection{Движение по окружности}

\subsubsection{Угловая скорость}

\textbf{Угловая скорость} ($\boldsymbol{\omega}$) — это векторная величина, характеризующая быстроту и направление вращения тела.
\[
\boldsymbol{\omega} = \frac{d\boldsymbol{\varphi}}{dt}
\]

Вектор $\boldsymbol{\omega}$ направлен вдоль оси вращения по \textit{правилу «буравчика»} (если вращать буравчик в сторону движения тела, поступательное движение винта укажет направление $\boldsymbol{\omega}$).

\subsubsection{Линейная скорость}

Линейная скорость точки $\mathbf{v}$, находящейся на радиус-векторе $\mathbf{r}$ от центра вращения, определяется векторным произведением:
\[
\mathbf{v} = [\boldsymbol{\omega}, \mathbf{r}] = \boldsymbol{\omega} \times \mathbf{r}
\]

\textbf{Доказательство.} Рассмотрим вращение точки в плоскости $XY$ вокруг оси $Z$ ($n_z$ --- единичный вектор).
Тогда векторы имеют вид:
\[
\boldsymbol{\omega} = \{0, 0, \omega\}, \quad \mathbf{r} = \{x, y, 0\}
\]

Координаты точки, движущейся по окружности:
\[
x = R \cos \varphi, \quad y = R \sin \varphi
\]

Дифференцируя по времени ($dt$), найдем проекции линейной скорости:
\[\begin{gathered}
v_x = \dot{x} = -R \sin \varphi \cdot \dot{\varphi} = -\omega y\\
v_y = \dot{y} = R \cos \varphi \cdot \dot{\varphi} = \omega x
\end{gathered}
\]

С другой стороны, вычислим векторное произведение $[\boldsymbol{\omega}, \mathbf{r}]$:
\[
[\boldsymbol{\omega}, \mathbf{r}] = 
\begin{vmatrix} 
\mathbf{i} & \mathbf{j} & \mathbf{k} \\ 
0 & 0 & \omega \\ 
x & y & 0 
\end{vmatrix} = 
\mathbf{i}(0 - \omega y) - \mathbf{j}(0 - \omega x) + \mathbf{k}(0) = \{-\omega y, \omega x, 0\}
\]

Совпадение результатов доказывает справедливость формулы. \( \square \)

\subsubsection{Зависимости $x, v_x, a_x$ от $\omega$}
Рассмотрим проекцию равномерного движения точки по окружности радиуса $R$ на ось $X$. Это движение является \textit{гармоническим колебанием}.
Если угол поворота $\varphi = \omega t$, то:

\begin{tblr}{
    colspec = {X[c]X[l]},
    vlines,
    hlines,
    row{1} = {font=\bfseries, mode=text},
    row{odd} = {bg=gray!10},
    cells = {m},
    rowsep = 6pt
}
Величина & Описание \\
$\displaystyle x(t) = R \cos(\omega t)$ & \textit{Координата} материальной точки, движущейся по окружности \\
$\displaystyle v_x(t) = \frac{dx}{dt} = -R \omega \sin(\omega t)$ & \textit{Проекция скорости} на ось $x$; модуль скорости $v = \omega R$ \\
$\displaystyle a_x(t) = \frac{dv_x}{dt} = -R \omega^2 \cos(\omega t)$ & \textit{Проекция ускорения} на ось $x$; модуль ускорения $a = \omega^2 R$ \\
$\displaystyle a = \frac{v^2}{R}$ & \textit{Связь ускорения и скорости}: центростремительное ускорение
\end{tblr}

\subsection{Момент инерции}

\subsubsection{Определение}

\textbf{Момент инерции} — скалярная физическая величина, мера инертности тела при вращательном движении.

\begin{center}
\begin{tblr}{
  colspec = {XX},
  vlines,
  hlines,
  row{1} = {font=\bfseries},
  row{odd} = {bg=gray!10},
  cells = {mode=math,c},
  rowsep = 6pt
}
\text{Поступательное движение} & \text{Вращательное движение} \\
\vec{r} & \alpha \\
\dfrac{\diff\vec{r}}{\diff t} = \vec{v}\quad\left[\dfrac{\text{м}}{\text{с}}\right] & \dfrac{\diff\vec{\alpha}}{\diff t} = \vec{\omega}\quad\left[\dfrac{\text{рад}}{\text{с}}\right] \\
\begin{array}{c} \vec{f} \\ \text{сила} \end{array} & \begin{array}{c} \vec{M}=[\vec{r}\times\vec{f}] \\ \text{момент силы} \end{array} \\
\begin{array}{c} m \\ \text{масса} \end{array} & \begin{array}{c} I \\ \text{момент инерции} \end{array} \\
\begin{array}{c} \vec{p} = m \vec{v} \\ \text{импульс} \end{array} & \begin{array}{c} \vec{L} = I \vec{\omega} \\ \text{момент импульса} \end{array} \\
\dfrac{\diff \vec{p}}{\diff t} = \vec{f} & \dfrac{\diff \vec{L}}{\diff t} = \vec{M} \\
\end{tblr}
\end{center}

\begin{center}
\begin{tblr}{
  colspec = {X[1.7,c]X[3,c]},
  vlines,
  hlines,
  row{1} = {font=\bfseries},
  row{odd} = {bg=gray!10},
  cells = {mode=math,c,m},
  rowsep = 8pt
}
\text{Тип системы} & \text{Формула момента инерции} \\
\text{Точечные массы} & I = \sum_{i} m_i r_i^2 = m_1r_1^2 + m_2r_2^2 + \cdots + m_nr_n^2 \\
\text{Непрерывное тело} & I = \int r^2  dm = \int_V \rho(\vec{r}) r^2  dV \\
\end{tblr}
\end{center}

\begin{center}
\begin{tblr}{
  colspec = {lcc},
  vlines,
  hlines,
  row{1} = {font=\bfseries,mode=text},
  row{odd} = {bg=gray!10},
  cells = {mode=math,c,m},
  rowsep = 4pt
}
\text{Тело} & \text{Ось вращения} & \text{Момент инерции} \\
\text{Стержень} & \text{Через центр} & \dfrac{1}{12}mL^2 \\
\text{Стержень} & \text{Через конец} & \dfrac{1}{3}mL^2 \\
\text{Кольцо} & \text{Через центр} & mR^2 \\
\text{Диск} & \text{Через центр} & \dfrac{1}{2}mR^2 \\
\text{Шар} & \text{Через центр} & \dfrac{2}{5}mR^2 \\
\text{Сфера} & \text{Через центр} & \dfrac{2}{3}mR^2 \\
\text{Цилиндр} & \text{Ось симметрии} & \dfrac{1}{2}mR^2 \\
\text{Пластина } a\times b & \text{Через центр} & \dfrac{1}{12}m(a^2 + b^2) \\
\end{tblr}
\end{center}

\subsubsection{Абсолютно твёрдое тело}

\textbf{Абсолютно твёрдым телом} называется система материальных точек, расстояние между которыми всегда остаётся постоянным:
\[
|\mathbf{r}_i - \mathbf{r}_j| = \text{const} \quad \text{для любых } i, j
\]

Ключевые свойства:
\begin{itemize}
    \item Тело не обладает внутренней потенциальной энергией упругой деформации ($\Delta U_{деф} = 0$).
    \item Движение тела полностью описывается движением его центра масс и вращением вокруг него.
    \item Работа внутренних сил в абсолютно твёрдом теле всегда равна нулю.
\end{itemize}

\subsubsection{Вывод момента импульса}

Рассмотрим абсолютно твердое тело, вращающееся вокруг неподвижной точки $O$ с угловой скоростью $\boldsymbol{\omega}$.

Момент импульса системы материальных точек определяется как сумма моментов импульса каждой точки:
\[
\mathbf{L} = \sum_i [\mathbf{r}_i, \mathbf{p}_i] = \sum_i m_i [\mathbf{r}_i, \mathbf{v}_i]
\]

Так как тело вращается как единое целое, линейная скорость каждой точки $\mathbf{v}_i$ выражается через угловую:
\[
\mathbf{v}_i = [\boldsymbol{\omega}, \mathbf{r}_i]
\]

Используем формулу двойного векторного произведения:
\[
[\mathbf{a}, [\mathbf{b}, \mathbf{c}]] = \mathbf{b}(\mathbf{a} \cdot \mathbf{c}) - \mathbf{c}(\mathbf{a} \cdot \mathbf{b})
\]
В нашем случае: $\mathbf{a} = \mathbf{r}_i$, $\mathbf{b} = \boldsymbol{\omega}$, $\mathbf{c} = \mathbf{r}_i$.

Раскрываем произведение для каждой точки:
\[
[\mathbf{r}_i, [\boldsymbol{\omega}, \mathbf{r}_i]] = \boldsymbol{\omega}(\mathbf{r}_i \cdot \mathbf{r}_i) - \mathbf{r}_i(\mathbf{r}_i \cdot \boldsymbol{\omega})
\]
Учитывая, что скалярное произведение вектора на самого себя дает квадрат его длины ($\mathbf{r}_i \cdot \mathbf{r}_i = r_i^2$):
\[
\mathbf{L} = \sum_i m_i \left( \boldsymbol{\omega} r_i^2 - \mathbf{r}_i (\mathbf{r}_i \cdot \boldsymbol{\omega}) \right)
\]

Пусть вращение происходит вокруг оси $Z$. Тогда $\boldsymbol{\omega} = (0, 0, \omega)$.
Радиус-вектор точки: $\mathbf{r}_i = (x_i, y_i, z_i)$.
Квадрат длины: $r_i^2 = x_i^2 + y_i^2 + z_i^2$.
Скалярное произведение: $(\mathbf{r}_i \cdot \boldsymbol{\omega}) = z_i \omega$.

Найдем проекцию момента импульса $L_z$ на ось $Z$. Для этого возьмем $z$-компоненту от общего уравнения:
\[
L_z = \sum_i m_i ( \omega r_i^2 - z_i (z_i \omega) ) = \omega \sum_i m_i ( r_i^2 - z_i^2 )
\]
Так как $r_i^2 - z_i^2 = (x_i^2 + y_i^2 + z_i^2) - z_i^2 = x_i^2 + y_i^2 = R_i^2$, где $R_i$ --- расстояние от точки до \textit{оси} вращения (а не до начала координат).

\[
L_z = \omega \sum_i m_i R_i^2
\]

Величина $\sum m_i R_i^2$ есть момент инерции относительно оси $Z$.
Таким образом, приходим к известной скалярной формуле:
\[\boxed{
L_z = I_z \omega
}\]

\subsection{Теорема Гюйгенса --- Штейнера}

Согласно \textbf{теореме Гюйгенса --- Штейнера}, момент инерции тела относительно произвольной оси равен сумме момента инерции относительно оси, проходящей через центр масс ($I_c$), и произведения массы тела на квадрат расстояния ($d$) между осями:

\begin{center}
\begin{tblr}{
  colspec = {XX},
  vlines,
  hlines,
  row{1} = {font=\bfseries},
  row{odd} = {bg=gray!10},
  row{2} = {valign=m},
  cells = {mode=math,c},
  rowsep = 6pt
}
\text{Формула} & \text{Схема} \\
I = I_{c} + md^2 
&
\begin{minipage}[c]{\linewidth}
\centering
\begin{tikzpicture}[scale=1.0]
    % Основные координаты
    \coordinate (B) at (1.5,0.8); % Центр масс
    \coordinate (A) at (3,1.5);   % Новая ось
    
    % Тело (произвольная фигура)
    \draw[fill=black!10, draw=black, thick] (B) circle (0.8);
    
    % Центр масс
    \filldraw[black] (B) circle (2pt) node[below right] {Ц.М.};
    
    % Новая ось
    \filldraw[black] (A) circle (2pt) node[above right] {Новая ось};
    
    % Расстояние d
    \draw[dashed, thick] (B) -- (A) node[midway, above] {$d$};
    
    % Обозначения моментов инерции
    \node at (1.5,2.2) {$I_{c}$};
    \node at (3.5,2.2) {$I$};
    
    % Стрелки к осям
    \draw[->, thick] (1.8,2) to[out=0,in=90] (B);
    \draw[->, thick] (3.2,2) to[out=180,in=90] (A);
    
    % Масса
    \node at (1,1.8) {$m$};
\end{tikzpicture}
\end{minipage}
\\
\end{tblr}
\end{center}